\documentclass{beamer}
\usepackage{mathtools}
\DeclarePairedDelimiter{\norm}{\lVert}{\rVert}

\usecolortheme{beaver}
\setbeamertemplate{blocks}[rounded=true, shadow=true]
\setbeamertemplate{footline}[page number]
%
\usepackage[utf8]{inputenc}
\usepackage[english,russian]{babel}
\usepackage{amssymb,amsfonts,amsmath,mathtext}
\usepackage{subfig}
\usepackage[all]{xy} % xy package for diagrams
\usepackage{array}
\usepackage{multicol} % many columns in slide
\usepackage{hyperref} % urls
\usepackage{hhline} %tables
\usepackage{comment} %comments
\usepackage{adjustbox}
\usepackage{multirow}
\usepackage{booktabs}
\newcommand{\blambda}{\boldsymbol{\lambda}}

\usepackage{amsmath,amssymb}

% Определения операторов (чтобы названия функций выглядели как операторы)
\DeclareMathOperator{\RMS}{RMS}
\DeclareMathOperator{\TB}{TB}
\DeclareMathOperator{\MSA}{MSA}
\DeclareMathOperator{\FFN}{FFN}
\newcommand{\Lmsa}{L_{\mathrm{MSA}}} % Было Lrms
\newcommand{\Lffn}{L_{\mathrm{FFN}}}

\DeclareMathOperator*{\argmin}{arg\,min}
\DeclareMathOperator*{\argmax}{arg\,max}

% Your figures are here:
\graphicspath{ {../figures/} }

%----------------------------------------------------------------------------------------------------------

\title[Научно-исследовательская работа]
{Научно-исследовательская работа}

\author[И.\,Д.~Степанов]{
    Степанов Илья Дмитриевич\\
    {\small
    Научный руководитель: к.\,ф.-м.\,н.\ А.\,В.~Грабовой\\
    }
}

\institute[МФТИ (НИУ)]{
    Кафедра интеллектуальных систем, ФПМИ МФТИ\\
    Московский физико-технический институт (национальный исследовательский университет)\\
}

\date{2026}



%----------------------------------------------------------------------------------------------------------
\begin{document}
%----------------------------------------------------------------------------------------------------------

\begin{frame}

    \thispagestyle{empty}
    \maketitle

\end{frame}

%-----------------------------------------------------------------------------------------------------

 \begin{frame}{Анализ Якобиана реккурентного трансформерного блока}


{\small
\textbf{Задача.} 
Провести анализ устойчивости итеративного трансформерного блока (MSA + FFN + LayerNorm) на основе исследования его якобиана.

\vspace{5pt}

\textbf{Проблема.} 
Известные результаты не учитывают полный состав компонентов блока, не хватает оценок спектральной нормы якобиана. Это затрудняет разработку итеративных трансформерных архитектур.

\vspace{5pt}

\textbf{Цель.} 
Вывести и доказать верхнюю оценку спектральной нормы якобиана полного блока для рекуррентных архитектур. Показать, что ключевую роль в подавлении роста градиентов и в обеспечении ограниченности спектральной нормы при больших шагах обновления играет слой LayerNorm. Подтвердить теоретические выводы численными экспериментами и провести сравнение с нормализацией RMSNorm.
}

\end{frame}
%------------------------------------------------------------------------------------------

\begin{frame}{Анализ Якобиана реккурентного трансформерного блока}

\small{
Рассматривается динамическая система:
\begin{flalign*}
  Y^{(t+1)} &= F\bigl(Y^{(t)}\bigr) = \RMS\bigl(U(Y^{(t)})\bigr)\\[4pt]
  U(Y)      &= Y + \nu\bigl(C + \TB(Y)\bigr)\\[4pt]
  Z         &= Y + \MSA(Y)\\[4pt]
  {F}_{\mathrm{msa}} &= \RMS(Z)\\[4pt]
  {F}_{\mathrm{ffn}} &= \FFN\bigl({F}_{\mathrm{msa}}\bigr)\\[4pt]
  \TB(Y)    &= \RMS\bigl({F}_{\mathrm{ffn}} + {F}_{\mathrm{msa}}\bigr)
\end{flalign*}
}
\small{
\textbf{Теорема 1 (Василенко)}:
Пусть входы в слои нормализации удовлетворяют условию нижней оценки их норм по строкам $\norm{U[i,:]} \ge r_U$, $\norm{Z[i,:]} \ge r_{in}$ и $\norm{({F}_{\mathrm{ffn}} + {F}_{\mathrm{msa}})[i,:]} \ge r_{out}$ для всех $i$. Тогда спектральная норма якобиана удовлетворяет неравенству:
\begin{equation*}
\lVert J_F(Y)\rVert_2 \le 
    \frac{\gamma_{\max}}{r_U}
    \left( 
        1 + 
        \nu \frac{\gamma_{\max}^2}{r_{in} r_{out}} (1 + \lVert J_{\FFN} \rVert_2)(1 + \lVert J_{\MSA} \rVert_2)
    \right).
\end{equation*}
где $\gamma_{\max}$ — максимальный параметр масштабирования RMSNorm.
}

\end{frame}

\begin{frame}{Анализ Якобиана реккурентного трансформерного блока}

% \begin{equation*}
\textbf{Теорема 2 (Степанов)}:
Пусть выполнены условия \textbf{теоремы 1}, и дополнительно предполагается условие
\[
\min_i \bigl\|C[i,:] + \operatorname{TB}(Y)[i,:]\bigr\|_2 \ge \varepsilon > 0.
\]
Тогда спектральная норма якобиана остаётся ограниченной при $\nu\to\infty$, то есть
\[
\|J_F(Y)\|_2 = O(1).
\]
% \end{equation*}
\

\vspace{10pt}
\color{red}\textbf{Работа принята на конференцию AINL 2026}
\end{frame}

\begin{frame}{Анализ Якобиана блока с LayerNorm: предложение}
\small
Рассматривается рекуррентный блок с LayerNorm:
\begin{flalign*}
  Y^{(t+1)} &= F\bigl(Y^{(t)}\bigr) = \operatorname{LayerNorm}\!\bigl(U(Y^{(t)})\bigr) \\[4pt]
  U(Y)      &= Y + \nu\bigl(C + \operatorname{TB}(Y)\bigr) \\[4pt]
  Z         &= Y + \operatorname{MSA}(Y) \\[4pt]
  {F}_{\mathrm{msa}} &= \operatorname{LayerNorm}(Z) \\[4pt]
  {F}_{\mathrm{ffn}} &= \operatorname{FFN}\bigl({F}_{\mathrm{msa}}\bigr) \\[4pt]
  \operatorname{TB}(Y) &= \operatorname{LayerNorm}\bigl({F}_{\mathrm{ffn}} + {F}_{\mathrm{msa}}\bigr)
\end{flalign*}

\vspace{4pt}
\textbf{Предлагается:}
\begin{itemize}
  \item Получить верхнюю оценку спектральной нормы якобиана \(\|J_F(Y)\|_2\)
        по аналогии с результатом для RMSNorm (Теорема 1).
  \item Установить, сохраняется ли ограниченность якобиана при \(\nu\to\infty\)
        (аналог Теоремы 2).
\end{itemize}
\end{frame}



%----------------------------------------------------------------------------------------------------------
\end{document}
